\documentclass[a4paper,11pt]{article}
\usepackage[utf8]{inputenc}
\usepackage[T1]{fontenc}
\usepackage[french]{babel}
\usepackage[left=1cm,right=1cm,top=.5cm,bottom=0cm]{geometry}
\usepackage{enumitem}
\usepackage{hyperref}
\usepackage{xcolor}
\usepackage{titlesec}
\usepackage{fontawesome5}
\usepackage{lmodern}
\usepackage[normalem]{ulem}

\definecolor{primary}{RGB}{220, 30, 30}
\definecolor{secondary}{RGB}{80, 80, 80}
\hypersetup{colorlinks=true, linkcolor=primary, filecolor=primary, urlcolor=primary}
\titleformat{\section}{\large\bfseries\color{primary}\uppercase}{}{0em}{}[\titlerule]
\titlespacing{\section}{0pt}{8pt}{5pt}

\begin{document}

\begin{center}
    {\Large \textbf{Lo\"ic Aron MBASSI EWOLO}} \\ \vspace{4pt}
    \textbf{\large Chef de Projet IA} \\
    \textit{\small Disponible d\`es septembre 2026 pour alternance 12 mois} \\ \vspace{4pt}
    \small
    \faMapMarker\ Cergy, France (Mobile IDF) \quad
    \faPhone\ (+33) 7 59 52 33 98 \quad
    \faEnvelope\ \href{mailto:loicaron.mbassiewolo@ensea.fr}{loicaron.mbassiewolo@ensea.fr} \\ \vspace{2pt}
    \faLinkedin\ \href{https://www.linkedin.com/in/loic-mbassi/}{linkedin.com/in/loic-mbassi} \quad
    \faGithub\ \href{https://github.com/Nameless0l}{github.com/Nameless0l} \quad
    \faGlobe\ \href{https://mbassiloic.tech}{mbassiloic.tech}
\end{center}

\vspace{-6pt}

\noindent\small{La t\'el\'ecommunication g\'en\`ere des volumes massifs de donn\'ees r\'eseau que les techniques NLP, Machine Learning et LLMs permettent d\'esormais de valoriser en temps r\'eel. Chez SFR, piloter des projets IA depuis l'id\'eation jusqu'\`a la mise en production -- en combinant ma\^itrise technique Python/ML et coordination d'\'equipes pluridisciplinaires -- repr\'esente une mont\'ee en comp\'etence directement align\'ee avec mon profil d'ing\'enieur-chercheur.}

\section{Formation}
\begin{itemize}[leftmargin=*, label=\tiny$\bullet$, nosep]
    \item \textbf{ENSEA -- \'Ecole Nationale Sup\'erieure de l'\'Electronique et de ses Applications} \hfill \textit{2025 -- 2027} \\
    \small{Double dipl\^ome d'Ing\'enieur -- Syst\`emes Embarqu\'es, Signal, IA.}
    \item \textbf{\'Ecole Polytechnique de Yaound\'e (1\textsuperscript{\`ere} \'ecole d'ing\'enieur CEMAC)} \hfill \textit{2021 -- 2025} \\
    \small{Dipl\^ome d'Ing\'enieur de Conception (Master 2) : \textbf{Math\'ematiques Appliqu\'ees}, G\'enie Logiciel, Algorithmique.}
\end{itemize}

\section{Comp\'etences Techniques}
\begin{itemize}[leftmargin=*, nosep]
    \item \small{\textbf{IA / NLP / LLM :} Python, PyTorch, BERT, Transformers, LangChain, RAG, analyse de donn\'ees, data viz}
    \item \small{\textbf{Gestion de projet IA :} cadrage, KPIs, suivi POC, coordination technique-m\'etier, documentation}
    \item \small{\textbf{Backend \& Infrastructure :} Docker, API REST, Laravel, NestJS, Git, CI/CD}
\end{itemize}

\section{Exp\'eriences \& Projets}
\begin{itemize}[leftmargin=*, label={-}]

\item \textbf{Stage de Recherche -- INRIA Saclay (\'Equipe TRIBE, IP Paris)} \hfill \textit{Avril -- Ao\^ut 2026} \\
\textit{Plateau de Saclay, France} \hfill \textit{Python, NLP, BERT, pipeline ML}
\vspace{-2pt}
    \begin{itemize}[leftmargin=0.4cm, nosep, label=\tiny$\bullet$]
        \item \small{Pilotage d'un pipeline NLP de bout en bout : d\'efinition des sp\'ecifications, impl\'ementation BERT, \'evaluation et reporting des r\'esultats.}
        \item \small{Coordination avec l'\'equipe de recherche pour l'adaptation du mod\`ele aux contraintes RGPD et la communication des r\'esultats.}
    \end{itemize}
\vspace{3pt}

\item \textbf{Head of Projects Cell \& Lead AI Cell -- ENSEA} \hfill \textit{2024 -- 2026} \\
\textit{Club de 50+ membres, Cergy} \hfill \textit{Leadership, gestion de projets IA, partenariats}
\vspace{-2pt}
    \begin{itemize}[leftmargin=0.4cm, nosep, label=\tiny$\bullet$]
        \item \small{Direction de la cellule projets ENSEA : planification, suivi et livraison de 5+ projets techniques en parall\`ele avec des \'equipes de 3-8 personnes.}
        \item \small{Mise en place de partenariats avec Google DeepMind et organisations internationales IA (IndabaX, MILA).}
    \end{itemize}
\vspace{3pt}

\item \textbf{Laravel API Generator -- Open Source} \hfill \textit{2024 -- 2025} \\
\textit{+30 \'etoiles GitHub} \hfill \textit{Laravel, LLM, UML, automatisation}
\vspace{-2pt}
    \begin{itemize}[leftmargin=0.4cm, nosep, label=\tiny$\bullet$]
        \item \small{Conception et livraison d'un outil open source de g\'en\'eration automatique d'APIs via LLM, avec gestion de roadmap, documentation et support communautaire.}
    \end{itemize}
\vspace{3pt}

\item \textbf{Urban Waste Management -- 1\textsuperscript{er} Prix CITS} \hfill \textit{2024} \\
\textit{Comp\'etition nationale (75 \'equipes)} \hfill \textit{ML, IoT, Python, gestion de projet}
\vspace{-2pt}
    \begin{itemize}[leftmargin=0.4cm, nosep, label=\tiny$\bullet$]
        \item \small{Coordination d'une \'equipe pluridisciplinaire pour d\'evelopper, tester et pr\'esenter une solution IA/IoT compl\`ete en moins de 48 heures.}
    \end{itemize}

\end{itemize}

\section{Prix \& Distinctions}
\begin{itemize}[leftmargin=*, nosep, label=\tiny$\bullet$]
    \item \small{\textbf{1\textsuperscript{er} prix IndabaX 2025} ; \textbf{1\textsuperscript{er} prix Mountain Hub 2024} ; \textbf{1\textsuperscript{er} prix CITS National} (vs 75 \'equipes)}
\end{itemize}

\end{document}
